\documentclass[10pt,xcolor=pst,aspectratio=169]{beamer}

\usepackage{etex}

%\usetheme{Boadilla}
%\usecolortheme{wolverine}
\usecolortheme{dolphin}
%\setbeamercovered{transparent}
%\setbeamercolor{block body}{bg=yellow}

\addtobeamertemplate{navigation symbols}{}{%
	\usebeamerfont{footline}%
	\usebeamercolor[fg]{footline}%
	\hspace{1em}%
	\insertframenumber/\inserttotalframenumber
}

\usepackage[utf8]{inputenc}
\usepackage[english,russian]{babel}
\usepackage[OT1]{fontenc}
\usepackage{amsmath, bm}
\usepackage{amsfonts}
\usepackage{amssymb}
\usepackage{graphicx}
\usepackage{wrapfig}
\usepackage[3D]{movie15}
\usepackage{animate}
\usepackage{ragged2e}
\usepackage{listings}
\usepackage{color}
\usepackage{pst-all}

\usepackage{tikz}
\usetikzlibrary{
	mindmap,
	arrows, % стрелки
	shapes.misc, % фигуры
	chains, % цепочки
	positioning, % позиционирование элементов
	scopes, % создание дополнительных веток
	shadows % тени
}

\graphicspath{{pic/}}

\author{\textbf{Губкин А.С.}}

\title[Численные методы в физике]{Численные методы в физике}

\logo{\includegraphics[width=0.1\linewidth]{LOGO_2.EPS}}

\institute[ТюмФ ИТПМ СО РАН]{Тюменский филиал Института теоретической и прикладной механики\\ им. С. А. Христиановича СО РАН, г. Тюмень}

%\date{6 октября 2015 г.}

\begin{document}

\lstset{ %
	language=[ANSI]C++,                 % выбор языка для подсветки (здесь это С++)
	keywordstyle=\color{blue},
	commentstyle=\color{gray},
	basicstyle=\scriptsize,
%basicstyle=\small\sffamily, % размер и начертание шрифта для подсветки кода
	numbers=left,               % где поставить нумерацию строк (слева\справа)
	numberstyle=\tiny,           % размер шрифта для номеров строк
%stepnumber=1,                   % размер шага между двумя номерами строк
	numbersep=4pt,                % как далеко отстоят номера строк от подсвечиваемого кода
%backgroundcolor=\color{white}, % цвет фона подсветки - используем \usepackage{color}
	showspaces=false,            % показывать или нет пробелы специальными отступами
	showstringspaces=false,      % показывать или нет пробелы в строках
	showtabs=false,             % показывать или нет табуляцию в строках
	frame=single,              % рисовать рамку вокруг кода
%tabsize=2,                 % размер табуляции по умолчанию равен 2 пробелам
	captionpos=t,              % позиция заголовка вверху [t] или внизу [b] 
	breaklines=true,           % автоматически переносить строки (да\нет)
	breakatwhitespace=true, % переносить строки только если есть пробел
	escapeinside={\%*}{*)}   % если нужно добавить комментарии в коде
}

%SLIDE #1
\begin{frame}

	\transdissolve[duration=0.1]
	\titlepage

\end{frame}

%SLIDE #?
\begin{frame}{Уравнение Бюргерса и Баклея - Леверетта}

	\transdissolve[duration=0.1]
	\justifying
	\large

	Будем рассматривать уравнение Бюргерса и Баклея - Леверетта. Соответствующи потоки:

	\[
		f(u) = \frac{1}{2} u^{2},
	\]

	\[
		f(u) = \frac{u^{2}}{u^{2} + \frac{1}{4}(1 - u^{2})^2}.
	\]

\end{frame}

%SLIDE #?
\begin{frame}{Уравнение Бюргерса и Баклея - Леверетта}

	\transdissolve[duration=0.1]
	\justifying
	\large

	Уравнение Бюргерса в консервативной форме:

	\[
		u_{t} + \left( \frac{1}{2} u^{2} \right)_{x} = 0.
	\]

	Уравнение Баклея - Леверетта в консервативной форме:

	\[
		u_{t} + \left(\frac{u^{2}}{u^{2} + \frac{1}{4}(1 - u^{2})^2} \right)_{x} = 0.
	\]

\end{frame}

%SLIDE #?
\begin{frame}{Пояснение к задаче Баклея - Леверетта}

	\transdissolve[duration=0.1]
	\justifying
	\large

	Уравнение Дарси для каждой фазы в дифференциальной форме:

	\[
		w_{1} = \frac{Q_{1}}{S} = - k \frac{k_{1}(s)}{\mu_{1}} p_{x},
	\]

	\[
		w_{2} = \frac{Q_{2}}{S} = - k \frac{k_{2}(s)}{\mu_{2}} p_{x}.
	\]

	Предполагается, что $s = s_{1}$ (либо $s = s_{2}$).\\

	Для двухфазного течения $s_{1}$ и $s_{2}$ связаны очевидным соотношением

	\[
		s_{1} + s_{2} = 1.
	\]

\end{frame}

%SLIDE #?
\begin{frame}{Пояснение к задаче Баклея - Леверетта}

	\transdissolve[duration=0.1]
	\justifying
	\large

	Уравнение неразрывности для каждой фазы:

	\[
		m s_{1t} + w_{1x} = 0,
	\]

	\[
		m s_{2t} + w_{2x} = 0.
	\]

	Складывая уравнения неразрывности, получим:

	\[
		\left( w_{1} + w_{2} \right)_{x} = 0,
	\]

	откуда найдем интеграл:

	\[
		w_{1} + w_{2} = w(t) \: \mbox{или} \: Q_{1} + Q_{2} = Q(t).
	\]

\end{frame}

%SLIDE #?
\begin{frame}{Пояснение к задаче Баклея - Леверетта}

	\transdissolve[duration=0.1]
	\justifying
	\large

	Исключим градиент давления:

	\[
		\frac{w_{1}}{w_{2}} = \frac{k_{1}(s)}{\frac{\mu_{1}}{\mu_{2}} k_{2}(s)}.
	\]

	Используя соотношение $w_{1} + w_{2} = w(t)$, получим

	\[
		\frac{w_{1}}{w(t)} = \frac{k_{1}(s)}{k_{1}(s) + \frac{\mu_{1}}{\mu_{2}} k_{2}(s)} = f(s) \: \mbox{- функция Баклея-Леверетта}.
	\]

	Поэтому

	\[
		w_{1} = f(s) w(t) \: \mbox{и} \: w_{2} = (1 - f(s)) w(t).
	\]

\end{frame}

%SLIDE #?
\begin{frame}{Пояснение к задаче Баклея - Леверетта}

	\transdissolve[duration=0.1]
	\justifying
	\large

	Подставим последнее соотношение для $w_{1}$ в уравнение неразрывности для первой фазы, получим:

	\[
		m s_{t} + w(t) f(s)_{x} = 0.
	\]

	Для случая, когда $w(t)/m = 1$

	\[
		s_{t} + f(s)_{x} = 0.
	\]

\end{frame}

%SLIDE #?
\begin{frame}{Уравнение Бюргерса}

	\transdissolve[duration=0.1]
	\justifying
	\large

	Рассмотрим аналитическое решение уравнения Бюргерса для различных начальных условий:

	\[
		\begin{split}
			&u_{t} + \left( \frac{1}{2} u^{2} \right)_{x} = 0, \\
			&f(u) = \frac{1}{2} u^{2}, \\
			&a(u) = u, \\
			&\frac{d x}{d t} = a(u), \\
			&x(t) = u_{0} (t - t_{0}) + x_{0}.
		\end{split}
	\]

\end{frame}

%SLIDE #?
\begin{frame}{Начальное условие №1}

	\transdissolve[duration=0.1]
	\justifying
	\large

	\[
		u(x,0) =
		\begin{cases}
			&\frac{1}{2} \qquad x = 0,\\
			&0 \qquad x \geq 0.
		\end{cases}
	\]

	\begin{center}
		\psset{xunit=4cm,yunit=4cm,algebraic=true}
		\begin{pspicture}(0,0)(2,1)
			\psgrid[griddots=20, gridwidth=0pt, gridcolor=gray, gridlabels=0pt, subgriddiv=5, subgriddots=20, subgridcolor=gray](0,0)(0,0)(2,1)
			\psaxes[Dx=0.2, Dy=0.2, subticks=2, labelFontSize=\scriptscriptstyle]{-}(0,0)(0,0)(2,1)[$x$,0][$u$,90]
			\psline[linewidth=2pt, linecolor=blue]{*}(0,0.5)(0,0.5)
			\uput[0](0,0.5){$u = \frac{1}{2}$}
			\psplot[linewidth=2pt, linecolor=blue, yMaxValue=1, yMinValue=0]{0}{2} {0}
		\end{pspicture}
	\end{center}

\end{frame}

%SLIDE #?
\begin{frame}{Поток уравнения Бюргерса и начальное условие}

	\transdissolve[duration=0.1]
	\justifying
	\large

	\begin{center}
		\psset{xunit=4cm,yunit=4cm,algebraic=true}
		\begin{pspicture}(-1,0)(1,1)
			\psgrid[griddots=20, gridwidth=0pt, gridcolor=gray, gridlabels=0pt, subgriddiv=2, subgriddots=20, subgridcolor=gray](0,0)(-1,0)(1,1)
			\psaxes[Dx=0.5, Dy=0.5, subticks=2, labelFontSize=\scriptscriptstyle]{->}(0,0)(-1,0)(1,1)[$u$,0][$f(u)$,90]
			\psplot[linewidth=2pt, linecolor=blue, yMaxValue=1, yMinValue=0]{-1}{1} {x^2}
			\psline[linewidth=2pt, linecolor=red]{o->}(0,0)(0.5,0.25)
			\uput[-45](0,0){$u = 0$}
			\uput[0](0.5,0.25){$u = \frac{1}{2}$}
			\uput[0]{26.56}(0.01,0.1){разрыв} 
		\end{pspicture}
	\end{center}

\end{frame}

%SLIDE #?
\begin{frame}{Харатеристики}

	\transdissolve[duration=0.1]
	\justifying
	\large

	Характеристики будут иметь вид:

	\[
		\frac{d x}{d t} = u \Rightarrow 
			\begin{cases}
				&u = u_0 =
				\begin{cases}
					&\frac{1}{2} \mbox{ для } t > 0 \: x > 0, \\
					&0 \mbox{ для } x > 0 \: t > 0,
				\end{cases} \\
				&x(t) = u_{0} t + x_{0} \mbox{ - начинаются на оси } \: x, \\
				&x(t) = u_{0} (t - t_{0}) \mbox{ - начинаются на оси } \: t.
			\end{cases}
	\]
	\textbf{До того момента пока они не пересекутся!}

\end{frame}

%SLIDE #?
\begin{frame}{Харатеристики}

	\transdissolve[duration=0.1]
	\justifying
	\large

	Интегрируя предыдущее уравнение, получим:

	\[
		x =
		\begin{cases}
			&\frac{1}{2} (t - t_{0}) \mbox{ для } t > 0 \: x > 0, \\
			&x_{0} \mbox{ для } x > 0 \: t > 0. 
		\end{cases}
	\]

\end{frame}

%SLIDE #?
\begin{frame}{Скорость распространения скачка}

	\transdissolve[duration=0.1]
	\justifying
	\large

	Раcчитаем скорость распространения скачка с помощью соотношения Ренкина - Гюгонио:

	\[
		S = \frac{[f]}{[u]} = \frac{f(u_{L}) - f(u_{R})}{u_{L} - u_{R}} = \frac{\frac{1}{2} \left( \frac{1}{2} \right)^{2} - \frac{1}{2} \left( 0 \right)^{2}}{\frac{1}{2} - 0} = \frac{1}{4}.
	\]

\end{frame}

%SLIDE #?
\begin{frame}{Аналитическое решение}

	\transdissolve[duration=0.1]
	\justifying
	\large

	Таким образом решение будет иметь следующий вид:

	\[
		u =
		\begin{cases}
			&\frac{1}{2}, \qquad t > 0, \qquad x > \frac{1}{4} t, \\
			&0, \qquad t > 0, \qquad x < \frac{1}{4} t.
		\end{cases}
	\]

	\begin{center}
		\psset{xunit=2cm,yunit=0.25cm,algebraic=true}
		\begin{pspicture}(0,0)(5,17)
			\psgrid[griddots=20, gridwidth=0pt, gridcolor=gray, gridlabels=0pt, subgriddiv=2, subgriddots=20, subgridcolor=gray](0,0)(0,0)(5,17)
			\psaxes[Dx=1, Dy=1, subticks=2, labelFontSize=\scriptscriptstyle]{->}(0,0)(0,0)(5,17)[$x$,0][$t$,90]

			\psline[linewidth=2pt, linecolor=green]{o->}(0,2)(1,4)
			\psline[linewidth=2pt, linecolor=green]{o->}(0,4)(2,8)
			\psline[linewidth=2pt, linecolor=green]{o->}(0,6)(3,12)
			\psline[linewidth=2pt, linecolor=green]{o->}(0,8)(4,16)

			\psplot[linewidth=2pt, linecolor=blue, yMaxValue=17, yMinValue=0]{0}{4} {4*x}

			\psline[linewidth=2pt, linecolor=red]{o->}(1,0)(1,4)
			\psline[linewidth=2pt, linecolor=red]{o->}(2,0)(2,8)
			\psline[linewidth=2pt, linecolor=red]{o->}(3,0)(3,12)
			\psline[linewidth=2pt, linecolor=red]{o->}(4,0)(4,16)
			\uput[0](0.5,15){$x = \frac{1}{2} (t - t_{0})$} 
			\uput[0](4,16){$x = \frac{1}{4} t$} 
			\uput[45]{90}(1,0){$x = 1$} 
			\uput[45]{90}(2,0){$x = 2$} 
			\uput[45]{90}(3,0){$x = 3$} 
			\uput[45]{90}(4,0){$x = 4$} 
		\end{pspicture}
	\end{center}

\end{frame}

%SLIDE #?
\begin{frame}{Аналитическое решение}

	\transdissolve[duration=0.1]
	\justifying
	\large

	\begin{center}
		\psset{xunit=1.2cm,yunit=5cm,algebraic=true}
		\begin{pspicture}(0,0)(8,1)
			\psgrid[griddots=20, gridwidth=0pt, gridcolor=gray, gridlabels=0pt, subgriddiv=2, subgriddots=20, subgridcolor=gray](0,0)(0,0)(8,1)
			\psaxes[Dx=1, Dy=0.5, subticks=2, labelFontSize=\scriptscriptstyle]{->}(0,0)(0,0)(8,1)[$x$,0][$u$,90]

			\psplot[linewidth=2pt, linecolor=blue, yMaxValue=1, yMinValue=0]{0}{7} {0.5}
			\psline[linewidth=2pt, linecolor=blue]{-}(1,0)(1,0.5)
			\psline[linewidth=2pt, linecolor=blue]{-}(2,0)(2,0.5)
			\psline[linewidth=2pt, linecolor=blue]{-}(3,0)(3,0.5)
			\psline[linewidth=2pt, linecolor=blue]{-}(4,0)(4,0.5)
			\psline[linewidth=2pt, linecolor=blue]{-}(5,0)(5,0.5)
			\psline[linewidth=2pt, linecolor=blue]{-}(6,0)(6,0.5)
			\psline[linewidth=2pt, linecolor=blue]{-}(7,0)(7,0.5)
			\uput[135]{90}(1,0){$t = 4$}
			\uput[135]{90}(2,0){$t = 8$}
			\uput[135]{90}(3,0){$t = 12$}
			\uput[135]{90}(4,0){$t = 16$}
			\uput[135]{90}(5,0){$t = 20$}
			\uput[135]{90}(6,0){$t = 24$}
			\uput[135]{90}(7,0){$t = 28$}
		\end{pspicture}
	\end{center}

\end{frame}

%SLIDE #?
\begin{frame}{Начальное условие №2}

	\transdissolve[duration=0.1]
	\justifying
	\large

	\[
		u(x,0) =
			\begin{cases}
				&0 \qquad x < -1,\\
				&\frac{1}{2} \qquad -1 < x < 0,\\
				&0 \qquad x > 0.
			\end{cases}
	\]

	\begin{center}
		\psset{xunit=3cm,yunit=3cm,algebraic=true}
		\begin{pspicture}(-2,0)(1,1)
			\psgrid[griddots=20, gridwidth=0pt, gridcolor=gray, gridlabels=0pt, subgriddiv=2, subgriddots=20, subgridcolor=gray](0,0)(-2,0)(1,1)
			\psaxes[Dx=0.5, Dy=0.5, subticks=2, labelFontSize=\scriptscriptstyle]{-}(0,0)(-2,0)(1,1)[$x$,0][$u$,90]

			\psplot[linewidth=2pt, linecolor=blue, yMaxValue=1, yMinValue=0]{-2}{-1} {0}
			\psline[linewidth=2pt, linecolor=blue]{-}(-1,0)(-1,0.5)
			\psplot[linewidth=2pt, linecolor=blue, yMaxValue=1, yMinValue=0]{-1}{0} {0.5}
			\psline[linewidth=2pt, linecolor=blue]{-}(0,0)(0,0.5)
			\psplot[linewidth=2pt, linecolor=blue, yMaxValue=1, yMinValue=0]{0}{1} {0}
			\uput[0](-0.7,0.7){$u = \frac{1}{2}$}
		\end{pspicture}
	\end{center}

\end{frame}

%SLIDE #?
\begin{frame}{Поток уравнения Бюргерса и начальное условие}

	\transdissolve[duration=0.1]
	\justifying
	\large

	\begin{center}
		\psset{xunit=4cm,yunit=4cm,algebraic=true}
		\begin{pspicture}(-1,0)(1,1)
			\psgrid[griddots=20, gridwidth=0pt, gridcolor=gray, gridlabels=0pt, subgriddiv=2, subgriddots=20, subgridcolor=gray](0,0)(-1,0)(1,1)
			\psaxes[Dx=0.5, Dy=0.5, subticks=2, labelFontSize=\scriptscriptstyle]{->}(0,0)(-1,0)(1,1)[$u$,0][$f(u)$,90]
			\psplot[linewidth=2pt, linecolor=blue, yMaxValue=1, yMinValue=0]{-1}{1} {x^2}
			\psline[linewidth=2pt, linecolor=red]{o->}(0,0)(0.5,0.25)
			\pscurve[linewidth=2pt, linecolor=red, yMaxValue=1, yMinValue=0]{->}(0.5,0.25)(0.375,0.140625)(0.25,0.0625)
			\uput[-45](0,0){$u = 0$}
			\uput[0](0.5,0.25){$u = \frac{1}{2}$}
			\uput[0]{26.56}(0.01,0.1){разрыв} 
			\uput[0]{0}(0.35,0.1){разрежение} 
		\end{pspicture}
	\end{center}
 
\end{frame}

%SLIDE #?
\begin{frame}{Харатеристики}

	\transdissolve[duration=0.1]
	\justifying
	\large

	Уравнение характеристик будет иметь вид:

	\[
		\frac{d x}{d t} = u_{0} =
			\begin{cases}
				&0 \qquad x_{0} < -1,\\
				&\frac{1}{2} \qquad -1 < x_{0} < 0,\\
				&0 \qquad x_{0} > 0,
			\end{cases}
	\]

	или

	\[
		\frac{d x}{d t} =
			\begin{cases}
				&0 \quad (u = 0) \qquad t > 0 \quad x < -1,\\
				&\frac{1}{2} \quad (u = \frac{1}{2}) \qquad t > 0 \quad -1 < x < 0,\\
				&0 \quad (u = 0) \qquad t > 0 \quad x > 0.
			\end{cases}
	\]

	пока они не пересекутся.

\end{frame}

%SLIDE #?
\begin{frame}{Харатеристики}

	\transdissolve[duration=0.1]
	\justifying
	\large

	Интегрируя предыдущее уравнение, получим:

	\[
		x =
			\begin{cases}
				&x_{0} \qquad t > 0 \quad x < -1,\\
				&\frac{1}{2}(t - t_{0})  \qquad t > 0 \quad -1 < x < 0,\\
				&x_{0} \qquad t > 0 \quad x > 0.
			\end{cases}
	\]

\end{frame}

%SLIDE #?
\begin{frame}{Скорость распространения скачка}

	\transdissolve[duration=0.1]
	\justifying
	\large

	Раcчитаем скорость распространения скачка с помощью соотношения Ренкина - Гюгонио:

	\[
		S = \frac{[f]}{[u]} = \frac{f(u_{L}) - f(u_{R})}{u_{L} - u_{R}} = \frac{ \frac{1}{2} \left( \frac{1}{2} \right)^{2} - \frac{1}{2}  \left( 0 \right)^{2}}{ \frac{1}{2} - 0} = \frac{1}{4}.
	\]

	Таким образом, разрыв движется из точки $x = 0$ со скоростью $S = \frac{1}{4}$.\\

	Траектория движения скачка $x_{s}(t)$  будет иметь вид:

	\[
		x_{s}(t) = 0 + \frac{1}{4} t = \frac{1}{4} t.
	\]

\end{frame}

%SLIDE #?
\begin{frame}{Веер волн разрежения}

	\transdissolve[duration=0.1]
	\justifying
	\large

	Чтобы удовлетворить энтропийному условию нужно заполнить "вакуум" между двумя семействами характеристик веером волн разрежения, который задается следующим уравнением:

	\[
		\frac{d x}{d t} = A = \frac{x + 1}{t},
	\]

	следовательно

	\[
		x + 1 = A t, \qquad 0 \leq A \leq \frac{1}{2}.
	\]

\end{frame}

%SLIDE #?
\begin{frame}{Веер волн разрежения}

	\transdissolve[duration=0.1]
	\justifying
	\large

	Волна разрежения с наклоном $A = \frac{1}{2}$ достигнет разрыв в момент времени:

	\[
		- 1 + \frac{1}{2} t = \frac{1}{4} t \Rightarrow t = 4.
	\]

	Левая часть этого уравнения представляет собой траекторию движения "хвоста" волны разрежения со значением наклона $u = \frac{1}{2}$ из точки $x = -1$, а правая - траектория движения разрыва.

\end{frame}

%SLIDE #?
\begin{frame}{Аналитическое решение для $t \leq 4$}

	\transdissolve[duration=0.1]
	\justifying
	\large

	$u = 0$ для $x < -1$ и справа от траектории движения разрыва $x_{s}(t)$.\\

	$u = \frac{1}{2}$ в области $\frac{1}{2} t - 1 < x < \frac{1}{4} t$.\\

	Таким образом, для промежутка времени $0 \leq t \leq 4$ решение будет иметь следующий вид

	\[
		u(x, t) =
			\begin{cases}
				&0 \qquad x < -1\\
				&\frac{x+1}{t} \qquad -1 < x < \frac{t}{2} - 1\\
				&\frac{1}{2} \qquad \frac{t}{2} - 1 < x < \frac{1}{4} t\\
				&0 \qquad \frac{1}{4} t > x
			\end{cases}.
	\]

\end{frame}

%SLIDE #?
\begin{frame}{Траектория скачка для $t \geq 4$}

	\transdissolve[duration=0.1]
	\justifying
	\large

	Рассмотрим траекторию движения скачка когда он пересечется с волной разрежения ($t \geq 4$). Для этого приравняем скорость движения скачка к средней скорости слева и справа от него:

	\[
		\frac{d x_{s}}{d t} = \frac{1}{2} \left( u_{L} + u_{R} \right) = \frac{1}{2} \left( \frac{x + 1}{t} + 0 \right).
	\]

	Таким образом, уравнение траектории скачка для $t \geq 4$:

	\[
		\frac{d x_{s}}{d t} = \frac{x + 1}{2 t}.
	\]

\end{frame}

%SLIDE #?
\begin{frame}{Траектория скачка для $t \geq 4$}

	\transdissolve[duration=0.1]
	\justifying
	\large

	Интегрируя последнее уравнение, получаем:

	\[
		\int \frac{d x_{s}}{x_{s} + 1} = \int \frac{d t}{2 t},
	\]

	\[
		\ln (x_{s} + 1) = \frac{1}{2} \ln (t) + \ln B,
	\]

	\[
		x_{s} + 1 = B t^{\frac{1}{2}}.
	\]

\end{frame}

%SLIDE #?
\begin{frame}{Траектория скачка для $t \geq 4$}

	\transdissolve[duration=0.1]
	\justifying
	\large

	Найдем константу $B$: мы знаем, что в момент времени $t = 4$ траектория движения скачка имеет координату $x = 1$, таким образм, $1 + 1 = B 4^{\frac{1}{2}} \Rightarrow B = 1$. \\

	Окончательно траектория движения скачка для $t \geq 4$, будет иметь вид:

	\[
		x_{s} = - 1 + t^{\frac{1}{2}}.
	\]

\end{frame}

%SLIDE #?
\begin{frame}{Аналитическое решение на характеристической плоскости}

	\transdissolve[duration=0.1]
	\justifying
	\large

	\begin{center}
		\psset{xunit=2cm,yunit=0.25cm,algebraic=true}
		\begin{pspicture}(-2,0)(3,16)
			\psgrid[griddots=20, gridwidth=0pt, gridcolor=gray, gridlabels=0pt, subgriddiv=2, subgriddots=20, subgridcolor=gray](0,0)(-2,0)(3,16)
			\psaxes[Dx=1, Dy=4, subticks=2, labelFontSize=\scriptscriptstyle]{->}(0,0)(-2,0)(3,16)[$x$,0][$t$,90]

			\psline[linewidth=0.5pt, linecolor=red]{o->}(-2,0)(-2,16)
			\psline[linewidth=0.5pt, linecolor=red]{o->}(-1.5,0)(-1.5,16)
			\psline[linewidth=0.5pt, linecolor=red]{o->}(-1,0)(-1,16)

			\psline[linewidth=0.5pt, linecolor=green]{o->}(-1,0)(-0.5,16)
			\psline[linewidth=0.5pt, linecolor=green]{o->}(-1,0)(0,16)
			\psline[linewidth=0.5pt, linecolor=green]{o->}(-1,0)(0.5,16)
			\psline[linewidth=0.5pt, linecolor=green]{o->}(-1,0)(1,16)
			\psline[linewidth=0.5pt, linecolor=green]{o->}(-1,0)(1.5,16)
			\psline[linewidth=0.5pt, linecolor=green]{o->}(-1,0)(2,16)
			\psline[linewidth=0.5pt, linecolor=green]{o->}(-1,0)(2.5,16)
			\psline[linewidth=0.5pt, linecolor=green]{o->}(-1,0)(1.5,6)
			\psline[linewidth=0.5pt, linecolor=green]{o->}(-1,0)(2,9)
			\psline[linewidth=0.5pt, linecolor=green]{o->}(-1,0)(2.5,12.25)
			\psline[linewidth=0.5pt, linecolor=green]{o->}(-1,0)(3,16)

			\psplot[linewidth=1pt, linecolor=blue, yMaxValue=4, yMinValue=0]{-1}{1} {2*(x + 1)}
			\psplot[linewidth=1pt, linecolor=blue, yMaxValue=4, yMinValue=0]{0}{1} {4*x}
			\psplot[linewidth=1pt, linecolor=blue, yMaxValue=16, yMinValue=0]{1}{3} {(x + 1)^2}

			\psline[linewidth=0.5pt, linecolor=red]{o->}(0.5,0)(0.5,2)
			\psline[linewidth=0.5pt, linecolor=red]{o->}(1,0)(1,4)
			\psline[linewidth=0.5pt, linecolor=red]{o->}(1.5,0)(1.5,6.25)
			\psline[linewidth=0.5pt, linecolor=red]{o->}(2,0)(2,9)
			\psline[linewidth=0.5pt, linecolor=red]{o->}(2.5,0)(2.5,12.25)
			\psline[linewidth=0.5pt, linecolor=red]{o->}(3,0)(3,16)

			\psline[linewidth=0.5pt, linecolor=black]{-}(-2,4)(3,4)
			\uput[0](-2,6){$t = 4$} 

			\uput[0](-0.5,13){веер волн разрежения} 
			\uput[0](-2,2){$\frac{d x}{d t} = 0$}
			\uput[0](2,2){$\frac{d x}{d t} = 0$}
			\uput[0]{45}(2,8){$x = - 1 + t^{\frac{1}{2}}$}
			\uput[0]{30}(0.25,0){$x = \frac{1}{4} t$}
			\uput[0]{15}(-0.9,2){$x = \frac{1}{2} t$}
		\end{pspicture}
	\end{center}

\end{frame}

%SLIDE #?
\begin{frame}{Уравнение Баклея - Леверетта}

	\transdissolve[duration=0.1]
	\justifying
	\large
	Рассмотрим уравнение Баклея - Леверетта в консервативной форме:
	\[
		u_{t} + \left(\frac{u^{2}}{u^{2} + \frac{1}{4}(1 - u^{2})^2} \right)_{x} = 0.
	\]
	Скорость распространения возмущений:
	\[
		\begin{split}
			&a(u) = f'(u) = \frac{(5 u^{2} - 2 u + 1) 8 u - 4 u^{2} (10 u - 2)}{(5 u^{2} - 2 u + 1)^{2}} =\\
			&= \frac{8 u - 8 u^{2}}{(5 u^{2} - 2 u + 1)^{2}}.
		\end{split}
	\]

\end{frame}

%SLIDE #?
\begin{frame}{Схема к задаче Баклея - Леверетта}

	\transdissolve[duration=0.1]
	\justifying
	\large

\end{frame}

%SLIDE #?
\begin{frame}{Поток уравнения Баклея - Леверетта}

	\transdissolve[duration=0.1]
	\justifying
	\large

	\begin{center}
		\psset{xunit=6cm,yunit=6cm,algebraic=true}
		\begin{pspicture}(0,0)(1,1)
			\psgrid[griddots=20, gridwidth=0pt, gridcolor=gray, gridlabels=0pt, subgriddiv=2, subgriddots=20, subgridcolor=gray](0,0)(0,0)(1,1)
			\psaxes[Dx=0.25, Dy=0.25, subticks=2, labelFontSize=\scriptscriptstyle]{->}(0,0)(0,0)(1,1)[$u$,0][$f(u)$,90]
			\psplot[linewidth=2pt, linecolor=blue, yMaxValue=1, yMinValue=0]{0}{1} {x^2/(x^2 + 0.25*(1 - x^2)^2)}
			\psplot[linewidth=2pt, linecolor=red, yMaxValue=1, yMinValue=0]{0}{1} {x^2}

			\uput[0]{50}(0.6,0.25){$f(u) = u^{2}$}
			\uput[90]{60}(0.15,0.25){$f(u) = \frac{u^{2}}{u^{2} + \frac{1}{4} (1 - u^{2})^{2}}$}
		\end{pspicture}
	\end{center}
 
\end{frame}

%SLIDE #?
\begin{frame}{Начальное условие №1}

	\transdissolve[duration=0.1]
	\justifying
	\large
	\[
		u(x,0) =
		\begin{cases}
			&\frac{1}{2} \qquad x = 0\\
			&0 \qquad x \geq 0
		\end{cases}
	\]

	\begin{center}
		\psset{xunit=4cm,yunit=4cm,algebraic=true}
		\begin{pspicture}(0,0)(2,1)
			\psgrid[griddots=20, gridwidth=0pt, gridcolor=gray, gridlabels=0pt, subgriddiv=5, subgriddots=20, subgridcolor=gray](0,0)(0,0)(2,1)
			\psaxes[Dx=0.2, Dy=0.2, subticks=2, labelFontSize=\scriptscriptstyle]{-}(0,0)(0,0)(2,1)[$x$,0][$u$,90]
			\psline[linewidth=2pt, linecolor=blue]{*}(0,0.5)(0,0.5)
			\uput[0](0,0.5){$u = \frac{1}{2}$}
			\psplot[linewidth=2pt, linecolor=blue, yMaxValue=1, yMinValue=0]{0}{2} {0}
\end{pspicture}
	\end{center}
 
\end{frame}

%SLIDE #?
\begin{frame}{Поток уравнения Баклея - Леверетта и начальное условие}

	\transdissolve[duration=0.1]
	\justifying
	\large

	\begin{center}
		\psset{xunit=6cm,yunit=6cm,algebraic=true}
		\begin{pspicture}(0,0)(1,1)
			\psgrid[griddots=20, gridwidth=0pt, gridcolor=gray, gridlabels=0pt, subgriddiv=2, subgriddots=20, subgridcolor=gray](0,0)(0,0)(1,1)
			\psaxes[Dx=0.25, Dy=0.25, subticks=2, labelFontSize=\scriptscriptstyle]{->}(0,0)(0,0)(1,1)[$u$,0][$f(u)$,90]
			\psplot[linewidth=2pt, linecolor=blue, yMaxValue=1, yMinValue=0]{0}{1} {x^2/(x^2 + 0.25*(1 - x^2)^2)}	
			\psline[linewidth=2pt, linecolor=red]{o->}(0,0)(0.5,0.64)

			\uput[-45](0,0){$u_{R} = 0$}
			\uput[0](0.5,0.64){$u_{L} = \frac{1}{2}$}
			\uput[0]{51}(0.1,0.25){разрыв} 
		\end{pspicture}
	\end{center}
 
\end{frame}

%SLIDE #?
\begin{frame}{Харатеристики}

	\transdissolve[duration=0.1]
	\justifying
	\large
	Уравнение характеристик будет иметь вид:
	\[
		\frac{d x}{d t} = a(u_{0}) = \frac{8 u_{0} - 8 u_{0}^{2}}{(5 u_{0}^{2} - 2 u_{0} + 1)^{2}}.
	\]
	Интегрируя это уравнение, получим характеристики:
	\[
		x = \left( \frac{8 u_{0} - 8 u_{0}^{2}}{(5 u_{0}^{2} - 2 u_{0} + 1)^{2}} \right) t + x_{0}.
	\]
	Подставляя начальное условие $u_{0} = \frac{1}{2}$, получим:
	\[
		x = \frac{32}{25} t + x_{0}.
	\]

\end{frame}

%SLIDE #?
\begin{frame}{Харатеристики}

	\transdissolve[duration=0.1]
	\justifying
	\large
	Окончательно, характеристики будут иметь вид:
	\[
		x =
		\begin{cases}
			&\frac{32}{25} t + x_{0} \mbox{ для } t > 0 \: x > 0 \mbox{ (берут начало на оси) } t, \\
			&x_{0} \mbox{ для } x > 0 \: t > 0 \mbox{ (берут начало на оси) } x. 
		\end{cases}
	\]

\end{frame}

%SLIDE #?
\begin{frame}{Скорость распространения разрыва}

	\transdissolve[duration=0.1]
	\justifying
	\large
	Раcчитаем скорость распространения разрыва с помощью соотношения Ренкина - Гюгонио:
	\[
		S = \frac{[f]}{[u]} = \frac{f(u_{L}) - f(u_{R})}{u_{L} - u_{R}} = \frac{ \left( \frac{\frac{1}{4}}{\frac{1}{4} + \frac{1}{4} \left(1 - \frac{1}{2} \right)^{2}} \right)^{2} - 0}{\frac{1}{2} - 0} = \frac{8}{5}.
	\]

\end{frame}

%SLIDE #?
\begin{frame}{Аналитическое решение на характеристической плоскости}

	\transdissolve[duration=0.1]
	\justifying
	\large
	Таким образом решение будет иметь следующий вид:
	\[
		u =
		\begin{cases}
			&\frac{1}{2}, \qquad t > 0, \qquad x > \frac{8}{5} t, \\
			&0, \qquad t > 0, \qquad x < \frac{8}{5} t.
		\end{cases}
	\]
	
	\begin{center}
		\psset{xunit=2cm,yunit=1cm,algebraic=true}
		\begin{pspicture}(0,0)(5,3)
			\psgrid[griddots=20, gridwidth=0pt, gridcolor=gray, gridlabels=0pt, subgriddiv=2, subgriddots=20, subgridcolor=gray](0,0)(5,3)
			\psaxes[Dx=1, Dy=1, subticks=2, labelFontSize=\scriptscriptstyle]{->}(0,0)(5,3)[$x$,0][$t$,90]

			\psline[linewidth=2pt, linecolor=green]{o->}(0,0.15625)(1,0.625)
			\psline[linewidth=2pt, linecolor=green]{o->}(0,0.3125)(2,1.25)
			\psline[linewidth=2pt, linecolor=green]{o->}(0,0.46875)(3,1.875)
			\psline[linewidth=2pt, linecolor=green]{o->}(0,0.625)(4,2.5)

			\psplot[linewidth=2pt, linecolor=blue, yMaxValue=3, yMinValue=0]{0}{5} {5/8*x}

			\psline[linewidth=2pt, linecolor=red]{o->}(1,0)(1,0.625)
			\psline[linewidth=2pt, linecolor=red]{o->}(2,0)(2,1.25)
			\psline[linewidth=2pt, linecolor=red]{o->}(3,0)(3,1.875)
			\psline[linewidth=2pt, linecolor=red]{o->}(4,0)(4,2.5)
			\uput[0](0.5,2.3){$x = \frac{32}{25} t + x_{0}$}
			\uput[0](4.2,2.3){$x = \frac{8}{5} t$}
			\uput[45]{90}(1,0){$x = 1$}
			\uput[45]{90}(2,0){$x = 2$}
			\uput[45]{90}(3,0){$x = 3$}
			\uput[45]{90}(4,0){$x = 4$}
		\end{pspicture}
	\end{center}

\end{frame}

%SLIDE #?
\begin{frame}{Аналитическое решение}

	\transdissolve[duration=0.1]
	\justifying
	\large	
	\begin{center}
		\psset{xunit=1.2cm,yunit=5cm,algebraic=true}
		\begin{pspicture}(0,0)(8,1)
			\psgrid[griddots=20, gridwidth=0pt, gridcolor=gray, gridlabels=0pt, subgriddiv=2, subgriddots=20, subgridcolor=gray](0,0)(0,0)(8,1)
			\psaxes[Dx=1, Dy=0.5, subticks=2, labelFontSize=\scriptscriptstyle]{->}(0,0)(0,0)(8,1)[$x$,0][$u$,90]

			\psplot[linewidth=2pt, linecolor=blue, yMaxValue=1, yMinValue=0]{0}{8} {0.5}
			\psline[linewidth=2pt, linecolor=blue]{-}(1.6,0)(1.6,0.5)
			\psline[linewidth=2pt, linecolor=blue]{-}(3.2,0)(3.2,0.5)
			\psline[linewidth=2pt, linecolor=blue]{-}(4.8,0)(4.8,0.5)
			\psline[linewidth=2pt, linecolor=blue]{-}(6.4,0)(6.4,0.5)
			\psline[linewidth=2pt, linecolor=blue]{-}(8,0)(8,0.5)
			\uput[135]{90}(1.6,0){$x = 1.6$}
			\uput[135]{90}(3.2,0){$x = 3.2$}
			\uput[135]{90}(4.8,0){$x = 4.8$}
			\uput[135]{90}(6.4,0){$x = 6.4$}
			\uput[135]{90}(8,0){$x = 8$}
			
			\uput[135]{90}(1.6,0.5){$t = 1$}
			\uput[135]{90}(3.2,0.5){$t = 2$}
			\uput[135]{90}(4.8,0.5){$t = 3$}
			\uput[135]{90}(6.4,0.5){$t = 4$}
			\uput[135]{90}(8,0.5){$t = 5$}

		\end{pspicture}
	\end{center}

\end{frame}

%SLIDE #?
\begin{frame}{Начальное условие №2}

	\transdissolve[duration=0.1]
	\justifying
	\large
	\[
		u(x,0) =
		\begin{cases}
			&\frac{3}{4} \qquad x = 0\\
			&0 \qquad x \geq 0
		\end{cases}
	\]

	\begin{center}
		\psset{xunit=4cm,yunit=4cm,algebraic=true}
		\begin{pspicture}(0,0)(2,1)
			\psgrid[griddots=20, gridwidth=0pt, gridcolor=gray, gridlabels=0pt, subgriddiv=5, subgriddots=20, subgridcolor=gray](0,0)(0,0)(2,1)
			\psaxes[Dx=0.2, Dy=0.2, subticks=2, labelFontSize=\scriptscriptstyle]{-}(0,0)(0,0)(2,1)[$x$,0][$u$,90]
			\psline[linewidth=2pt, linecolor=blue]{*}(0,0.75)(0,0.75)
			\uput[0](0,0.75){$u = \frac{3}{4}$}
			\psplot[linewidth=2pt, linecolor=blue, yMaxValue=1, yMinValue=0]{0}{2} {0}
\end{pspicture}
	\end{center}
 
\end{frame}

%SLIDE #?
\begin{frame}{Поток уравнения Баклея - Леверетта и начальное условие}

	\transdissolve[duration=0.1]
	\justifying
	\large

	\begin{center}
		\psset{xunit=6cm,yunit=6cm,algebraic=true}
		\begin{pspicture}(0,0)(1,1)
			\psgrid[griddots=20, gridwidth=0pt, gridcolor=gray, gridlabels=0pt, subgriddiv=2, subgriddots=20, subgridcolor=gray](0,0)(0,0)(1,1)
			\psaxes[Dx=0.25, Dy=0.25, subticks=2, labelFontSize=\scriptscriptstyle]{->}(0,0)(0,0)(1,1)[$u$,0][$f(u)$,90]
			\psplot[linewidth=2pt, linecolor=blue, yMaxValue=1, yMinValue=0]{0}{1} {x^2/(x^2 + 0.25*(1 - x^2)^2)}	
			\psline[linewidth=2pt, linecolor=red]{o->}(0,0)(0.75,0.9216)

			\uput[-45](0,0){$u_{R} = 0$}
			\uput[0](0.75,0.9216){$u_{L} = \frac{3}{4}$}
			\uput[0]{50}(0.1,0.25){разрыв}
			\uput[90](0.25,0.8){!?}

			\psline[linewidth=2pt, linecolor=red]{->}(0.25,0.8)(0.55,0.75)
		\end{pspicture}
	\end{center}
 
\end{frame}

%SLIDE #?
\begin{frame}{Поток уравнения Баклея - Леверетта и начальное условие}

	\transdissolve[duration=0.1]
	\justifying
	\large

	\begin{center}
		\psset{xunit=6cm,yunit=6cm,algebraic=true}
		\begin{pspicture}(0,0)(1,1)
			\psgrid[griddots=20, gridwidth=0pt, gridcolor=gray, gridlabels=0pt, subgriddiv=2, subgriddots=20, subgridcolor=gray](0,0)(0,0)(1,1)
			\psaxes[Dx=0.25, Dy=0.25, subticks=2, labelFontSize=\scriptscriptstyle]{->}(0,0)(0,0)(1,1)[$u$,0][$f(u)$,90]
			\psplot[linewidth=2pt, linecolor=blue, yMaxValue=1, yMinValue=0]{0}{1} {x^2/(x^2 + 0.25*(1 - x^2)^2)}	
			\psline[linewidth=2pt, linecolor=red]{o->}(0,0)(0.617403,0.799268753)
			\psplot[linewidth=2pt, linecolor=red, yMaxValue=1, yMinValue=0]{0.617403}{0.75} {x^2/(x^2 + 0.25*(1 - x^2)^2)}
			\psline[linewidth=2pt, linecolor=red]{o}(0.617403,0.799268753)(0.617403,0.799268753)
			\psline[linewidth=2pt, linecolor=red]{o}(0.75,0.9216)(0.75,0.9216)

			\uput[-45](0,0){$u_{R} = 0$}
			\uput[0](0.75,0.9216){$u_{L} = \frac{3}{4}$}
			\uput[0](0.617403,0.799268753){$u_{T}$}
			\uput[0]{52}(0.1,0.25){разрыв}
			\uput[90](0.25,0.9){разрежение}

			\psline[linewidth=2pt, linecolor=red]{->}(0.25,0.9)(0.65,0.85)
		\end{pspicture}
	\end{center}
 
\end{frame}

%SLIDE #?
\begin{frame}{Волна разрежения}

	\transdissolve[duration=0.1]
	\justifying
	\large
	В данном случае в решение будет присутствовать волна разрежения, поскольку нельзя соединить точки $(0,0)$ и $\left( \frac{3}{4},f \left( \frac{3}{4} \right) \right)$ одной линией и удовлетворить условию энтропии.\\
	Поэтому нужно найти точку $\left( u_{T}, f \left( u_{T} \right) \right)$ в которой наклон касательной к кривой $f(u)$ будет совпадать с наклоном прямой, соединящей точки $(0,0)$ и $\left( u_{T}, f \left( u_{T} \right) \right)$:
	\[
		f' \left( u_{T} \right) = \frac{f \left( u_{T} \right) - 0}{u_{T} - 0}.
	\]

\end{frame}

%SLIDE #?
\begin{frame}{Волна разрежения}

	\transdissolve[duration=0.1]
	\justifying
	\large
	Получим следующее нелинейное уравнение:
	\[
		\frac{8 u_{T} - 8 u_{T}^2}{\left( 5 u_{T}^{2} - 2 u_{T} + 1 \right)^{2}} = \frac{u_{T}}{u_{T}^{2} + \frac{1}{4} \left(1 - u_{T}^{2} \right)^{2}}.
	\]
	Его решение можно получить, например, с помощью метода Ньютона:
	\[
		\begin{split}
			&F \left( u_{T} \right) = \frac{8 u_{T} - 8 u_{T}^2}{\left( 5 u_{T}^{2} - 2 u_{T} + 1 \right)^{2}} - \frac{u_{T}}{u_{T}^{2} + \frac{1}{4} \left(1 - u_{T}^{2} \right)^{2}},\\
			& u_{T}^{(n + 1)} = u_{T}^{(n)} - \frac{ F \left( u_{T}^{(n)} \right)}{F' \left( u_{T}^{(n)} \right)} \Rightarrow u_{T} \approx 0.617403.
		\end{split}
	\]

\end{frame}

%SLIDE #?
\begin{frame}{Аналитическое решение на характеристической плоскости}

	\transdissolve[duration=0.1]
	\justifying
	\large

	\begin{center}
		\psset{xunit=2cm,yunit=2cm,algebraic=true}
		\begin{pspicture}(0,0)(3,3)
			\psgrid[griddots=20, gridwidth=0pt, gridcolor=gray, gridlabels=0pt, subgriddiv=2, subgriddots=20, subgridcolor=gray](0,0)(3,3)
			\psaxes[Dx=1, Dy=1, subticks=2, labelFontSize=\scriptscriptstyle]{->}(0,0)(3,3)[$x$,0][$t$,90]

			\psline[linewidth=2pt, linecolor=green]{o->}(0,0.5)(0.5,1)
			\psline[linewidth=2pt, linecolor=green]{o->}(0,1)(1,2)
			\psline[linewidth=2pt, linecolor=green]{o->}(0,1.5)(1.5,3)
			\psline[linewidth=2pt, linecolor=green]{o->}(0,2)(1,3)
			\psline[linewidth=2pt, linecolor=green]{o->}(0,2.5)(0.5,3)


			\psplot[linewidth=2pt, linecolor=blue, yMaxValue=3, yMinValue=0]{0}{3} {2*x}

			\psplot[linewidth=2pt, linecolor=yellow, yMaxValue=3, yMinValue=0]{0}{3} {1.7*x}
			\psplot[linewidth=2pt, linecolor=yellow, yMaxValue=3, yMinValue=0]{0}{3} {1.4*x}
			\psplot[linewidth=2pt, linecolor=yellow, yMaxValue=3, yMinValue=0]{0}{3} {1.1*x}
			\psplot[linewidth=2pt, linecolor=yellow, yMaxValue=3, yMinValue=0]{0}{3} {0.9*x}

			\psplot[linewidth=2pt, linecolor=blue, yMaxValue=3, yMinValue=0]{0}{3} {625/768*x}

			\psline[linewidth=2pt, linecolor=red]{o->}(1,0)(1,0.813802083)
			\psline[linewidth=2pt, linecolor=red]{o->}(2,0)(2,1.627604167)
			\psline[linewidth=2pt, linecolor=red]{o->}(3,0)(3,2.44140625)

			\uput[0]{50}(0.6,1){веер волн разрежения}
			\uput[0]{40}(0.5,0.2){траектория движения разрыва}
		\end{pspicture}
	\end{center}

\end{frame}

\end{document}